\chapter{\acs*{FIBnt}}

\section{Process optimisation}

Due to the destructive nature of the \ac{FIBnt} process, there are a lot of possible artefacts.

\begin{figure}[H]
    \centering

	\begin{subfigure}[t]{0.32\textwidth}
		\includegraphics[width=\textwidth]{methods/fib_artifacts_b.pdf}
		\caption{Redeposition in front of the \ac{ROI} and slight vertical shadowing.}
		\label{fig:fib_artifacts_b}
	\end{subfigure}
	\hfill
	\begin{subfigure}[t]{0.32\textwidth}
		\includegraphics[width=\textwidth]{methods/fib_artifacts_c.pdf}
		\caption{Charging effects in form of drift artifacts and redepositioning.}
		\label{fig:fib_artifacts_c}
	\end{subfigure}
	\hfill
	\begin{subfigure}[t]{0.32\textwidth}
		\includegraphics[width=\textwidth]{methods/fib_artifacts_a.pdf}
		\caption{Curtaining and charging effects.}
		\label{fig:fib_artifacts_a}
	\end{subfigure}
    \caption{Imaging issues of hydrated and embedded \ac{C3S} specimens during the \ac{FIBnt}-process within the bulk material.}
    \label{fig:fib_artifacts}
\end{figure}

As in the \ac{ArBIB} sectioning (see \zcref{sec:bib_pol_fixed}), the cutting process can introduce curtaining.
It is caused by the cutting process causing vertical line artifacts in the direction of the ion beam (\zcref{fig:fib_artifacts_a}).
As in images from the stationary \ac{ArBIB} process, these lines can be reduced in post processing using a directional \ac{FFT} based filter \cite{fitschen.2017,schankula.2018}.

Vertical shadowing is another typical artifact caused by a loss of signal of the back-scattered electrons due to the trench geometry (see \zcref{fig:fib_artifacts_b}).
This well known effect can be corrected in post processing \cite{holzer.2007, mac.2018, fager.2020}, if the phase contrast is good enough.
However, there are more severe uncorrectable artifacts caused by the in-place slicing process.
Since cementitious material is not conductive and the trenches are not coated throughout, charging effects can negatively influence the image quality (e.g. \zcref{fig:fib_artifacts_c} and \zcref{fig:fib_artifacts_a}) or cause uncorrectable image drift during the process.
The re-deposition of the removed material usually happens within the trenches, which does not cause any problems.
However, it may also appear on small structures in front of the \ac{ROI} or even on the specimen surface after longer imaging periods as shown in \zcref{fig:fib_artifacts_b} and \zcref{fig:fib_artifacts_c}.

To avoid or at least reduce issues caused by slicing within the bulk material, the \ac{ROI} can be completely cut free.
Then a lift-out process can be applied, which is similar to the preparation process of \ac{TEM} lamellas \cite{overwijk.1993, langford.2001a}.
To do this, a needle is inserted and welded to the \ac{ROI}-cube using the platinum coating process.
Afterwards, the cube is complete cut loose from the bulk material, and a copper comb is inserted to the SEM.
Finally, the cube is transferred and then welded to a finger of the comb as shown in \zcref{fig:fib_liftout}.

\begin{figure}[h!]
    \centering

	\includegraphics[width=\textwidth]{methods/fib_liftout/image095_colored.pdf}
    \caption{\qty{20}{\micro\meter} wide specimen ({\color{green}A}).
	On the left is the needle welded to the specimen ({\color{blue}B}). They are navigated and welded to the copper comb ({\color{red}C}).
	The platinum vapour source ({\color{orange}D}) used for welding is visible on the right. \cite{kleiner.2021c}}
    \label{fig:fib_liftout}
\end{figure}

The so prepared and separated cube can now be serial cut as illustrated in \zcref{fig:fib_scheme} with significantly reduced shadowing and charging effects.

A detailed video of the \ac{FIBnt}-process, the required parameters, and potential sources of errors for cement-based materials was created and published \cite{kleiner.2025a}.

\section{Evaluation obstacles}

Due to the sensitive nature of the hydrate phases (see \zcref{sec:beam_damage}), the operator must pay particular attention when analysing at higher voltages as for \ac{EDX} or \ac{EBSD} samples.
As shown in \zcref{fig:edx_damage}, an electron beam at \qty{6}{\kilo\volt} and a long exposure time can cause damage in the specimen which could be misinterpreted as porosity.
The exposure time is the sum of the dwell time on each point and the frame count.
But especially the frame count increases the risk of specimen damage.
In the example, pores with a diameter of several \unit{\nano\meter} were induced, which can mostly be identified by comparing the damaged slice (\zcref{fig:edx_damage_a}) with a freshly cut surface (\zcref{fig:edx_damage_b}).
In this case, the damage was visible for $\approx$ \qty{350}{\nano\meter}

\begin{figure}[H]
    \centering

	\begin{subfigure}[t]{0.49\textwidth}
		\includegraphics[width=\textwidth]{methods/edx-damage/a.pdf}%
		\caption{A undamaged surface slice.}
		\label{fig:edx_damage_a}
	\end{subfigure}
	\hfill
	\begin{subfigure}[t]{0.49\textwidth}
		\includegraphics[width=\textwidth]{methods/edx-damage/b.pdf}
		\caption{Damaged surface after \ac{EDX}-measurements.}
		\label{fig:edx_damage_b}
	\end{subfigure}
    \caption{
        Effect of \ac{EDX} at \qty{6}{\kilo\volt} after multiple scans on a hydrated \ac{C3S} surface, prepared using \ac{FIB}.
        Sections of \ac{BSE}-images created during a \ac{FIBnt}-process, (\qty{2}{\kilo\volt}, \qty{2}{\pico\ampere}).
    } % BeamCurrent=2e-010
    \label{fig:edx_damage}
\end{figure}

Solutions to this problem are to reduce the scanning duration or the electron beam current, which subsequently will reduce the overall counts, to reduce the electron beam voltage, which will limit the elements which can be analysed or to use more sensitive detectors.
For the presented experiments, a even more reduced electron beam voltage was not applicable.
Therefore, a very sensitive, windowless \ac{EDX}-detector was used and combined with the counts of a second, less-sensitive \ac{EDX}-detector, which yields only one-fifth of the counts of the windowless detector.
The exposure time (\qty{60}{\micro\second}, 60 frames) and the current were reduced as low as possible, while leaving the beam voltage at \qty{10}{\kilo\volt}.

This still results in a very noisy, low count \ac{EDX} dataset.
However, it allows the distinction of the different main phases of the material.

\section{Postprocessing}

The so acquired datasets require much less post processing to reduce artefacts.
However, some problems still remain to be solved.
For example, the noisy \ac{EDX} dataset has to be denoised with an edge preserving denoising algorithm.
Since the signal-to noise ratio of the \ac{EDX} element mappings are very low, algorithms like median and \ac{NLM} are not suitable.
\ac{BM3D} denoising was found to be effective for the noise found in the \ac{EDX} slices and enables to reproducible quantitative analysis across elements.

The main challenge for the dataset is the precise alignment of the acquired images.
Drift during acquisition can cause a shift of the observed object between consecutive slices.
For a single image stack, this issue can be addressed effectively using the \ac{SIFT} algorithm as implemented in \textit{TrakEM2} \cite{cardona.2012}, which allows for automated correction of in-plane translations.

However, aligning the \ac{EDX} dataset to the corresponding \ac{BSE} dataset is more challenging.
Although the correspondence between \ac{BSE} and \ac{EDX} slices is known, the images differ in resolution and field of view.
If the image shift between the \ac{BSE} and \ac{EDX} images of the same slice is nearly identical, the alignment problem can be solved relatively easily.
%Otherwise, additional registration steps, such as the use of landmarks or scale-adjusted transformations, are required to achieve accurate overlay of the datasets.

Therefore, a python tool utilising \textit{napari} \cite{sofroniew.2025}, a graphical user interface to view and edit 2D and 3D image data, was developed.
This workflow aligns, visualizes, and analyses correlative \ac{BSE} and \ac{EDX} image stacks.
First, it loads the preprocessed in-plane translation data stored in the \textit{TrakEM2} project and the user has to provide the image scale and the slice thickness in \unit{\nano\meter\per\px}, if these information are not stored in the metadata of the files.
After denoising, landmark pairs between the first \ac{BSE} and \ac{EDX} image have to be provided by the user.
This is then used to estimate an isotropic scale-and-translation transform for a good stack-to-stack alignment.

Due to the geometry within the \ac{SEM}, the image stacks inherently have a systematically shift in the vertical direction \ac{SIFT}-algorithm.
This can be compensated using a user-defined angle.
This angle depends mostly on the device geometry and the angle at which the cube is welded to the copper comb.

{
	\color{red} % has to be rewritten
	Inter-slice intensity banding ($z$-banding) in the \ac{EDX} and brightness changes in the \ac{BSE} datasets are reduced via a normalization based on the mean brightness value of the overall stack.
	The registered 3D \ac{BSE} and \ac{EDX} volumes are displayed in \textit{napari} with physically meaningful voxel scales.
	The user then has to select an \ac{ROI}, which will be used for the further processing.
	Within the \ac{ROI}, 3D \ac{SLIC}O superpixels are computed on the denoised \ac{EDX} volume.
	Similar to the basic idea of Multi\ac{SLIC}, three main elements (\ce{Ca}, \ce{Si}, \ce{Al}) are selected and set as individual RGB color channels in an image, from which the superpixels are generated.
	This allows to use the standard implementations of \ac{SLIC} as found in \textit{OpenCV}.
	For each superpixel, mean \ac{EDX} intensities are aggregated for each element from the raw \ac{EDX}-data.
	These superpixel-wise vectors are then standardised and clustered using $k$-means in combination with \ac{UMAP}.
	This allows to identify material phases, yielding a 3D phase label volume that is directly overlaid in \textit{napari}.
}